\documentclass[10pt,a4paper]{article}

\usepackage[T1]{fontenc}
\usepackage[utf8]{inputenc}
\usepackage[ngerman]{babel}
\usepackage[a4paper,margin=14mm]{geometry}
\usepackage{lmodern}
\usepackage[table]{xcolor}
\usepackage{tabularx}
\usepackage{array}
\usepackage{microtype}
\usepackage[most]{tcolorbox}

\pagestyle{empty}
\setlength{\parindent}{0pt}
\setlength{\parskip}{0pt}
\renewcommand{\familydefault}{\sfdefault}
\renewcommand{\arraystretch}{1.16}

\definecolor{Navy}{HTML}{17324D}
\definecolor{Blue}{HTML}{3B6EA5}
\definecolor{Sky}{HTML}{EAF3FA}
\definecolor{Mint}{HTML}{E8F5EF}
\definecolor{Gold}{HTML}{F3B33D}
\definecolor{Ink}{HTML}{1E2935}
\definecolor{Muted}{HTML}{617080}
\definecolor{Line}{HTML}{D8E1E8}
\definecolor{CodeBg}{HTML}{F5F7F9}

\color{Ink}
\newcommand{\sectiontitle}[1]{%
  \vspace{2.3mm}{\color{Blue}\bfseries\large #1}\par
  \vspace{0.7mm}{\color{Blue}\rule{\linewidth}{0.7pt}}\par\vspace{1.2mm}}
\newcommand{\schema}[1]{\colorbox{CodeBg}{\parbox{\dimexpr\linewidth-2\fboxsep\relax}{\ttfamily\small #1}}}

\begin{document}

\colorbox{Navy}{\parbox{\dimexpr\linewidth-2\fboxsep\relax}{%
  \vspace{3mm}\color{white}\begin{tabularx}{\linewidth}{@{}Xr@{}}
  {\small\bfseries INFORMATIK · KLASSE 10} & \colorbox{Gold}{\color{Navy}\small\bfseries\strut LÖSUNGEN}\end{tabularx}
  \vspace{2.4mm}{\fontsize{18}{21}\selectfont\bfseries Aufgabe 8 -- 1:1- und n:m-Beziehungen}\par
  \vspace{0.8mm}{\large Sicherungsblatt}\vspace{2.8mm}}}

\vspace{2.5mm}
\colorbox{Sky}{\parbox{\dimexpr\linewidth-2\fboxsep\relax}{\vspace{1.6mm}
\textcolor{Blue}{\bfseries Ziel:} Du erkennst 1:1- und n:m-Beziehungen und setzt sie mit Fremdschlüsseln und Hilfstabellen um.\vspace{1.6mm}}}

\sectiontitle{1 · 1:1-Beziehung: Klasse und Klassenleiter}
{\small\texttt{klasse.id = 1} (10a) gehört zu \texttt{klassenleiter.klasse\_id = 1} (KOH). Ebenso: 10b -- SEI und 10c -- BRA. Jede Klasse hat genau einen Klassenleiter; jeder Klassenleiter leitet genau eine Klasse.}\par\vspace{1.1mm}
{\centering\bfseries\textcolor{Blue}{Klasse \quad 1 \;\rule{32mm}{0.7pt}\; 1 \quad Klassenleiter}\par}
\vspace{1.2mm}
\begin{minipage}[t]{0.48\linewidth}
\textcolor{Blue}{\bfseries Möglichkeit A}\par\vspace{0.6mm}
\schema{klasse (id: int, name: varchar(255))\par
klassenleiter (id: int, kuerzel: varchar(255),\par\quad klasse\_id[klasse]: int)}
\end{minipage}\hfill
\begin{minipage}[t]{0.48\linewidth}
\textcolor{Blue}{\bfseries Möglichkeit B}\par\vspace{0.6mm}
\schema{klasse (id: int, name: varchar(255),\par\quad klassenleiter\_id[klassenleiter]: int)\par
klassenleiter (id: int, kuerzel: varchar(255))}
\end{minipage}
\vspace{1.2mm}\par
{\small\textbf{Merke:} Bei 1:1 kann der Fremdschlüssel in einer der beiden Tabellen stehen. Ein Fremdschlüssel genügt; er wird nicht in beide Tabellen eingetragen.}\par

\sectiontitle{2 · n:m-Beziehung: Schüler und AG}
{\small Lina besucht Robotik und Theater; die Theater-AG besuchen Lina und Mia. Ein einzelner Fremdschlüssel in \texttt{schueler} oder \texttt{ag} könnte die mehreren Zuordnungen nicht speichern.}\par\vspace{0.8mm}
{\centering\bfseries\textcolor{Blue}{Schüler \quad n \;\rule{32mm}{0.7pt}\; m \quad AG}\par}
\vspace{1mm}
\begin{minipage}[t]{0.42\linewidth}
\textcolor{Blue}{\bfseries Hilfstabelle}\par\vspace{0.6mm}
\schema{teilnahme (\par\quad schueler\_id[schueler]: int\par\quad ag\_id[ag]: int\par)}
\end{minipage}\hfill
\begin{minipage}[t]{0.55\linewidth}
\textcolor{Blue}{\bfseries Beispiele aus teilnahme}\par\vspace{0.6mm}
{\small\ttfamily
\begin{tabular}{|c|c|l|}\hline
\rowcolor{Sky}schueler\_id & ag\_id & Zuordnung\\\hline
1 & 1 & Lina -- Robotik\\
1 & 2 & Lina -- Theater\\
3 & 2 & Mia -- Theater\\\hline
\end{tabular}}
\end{minipage}
\vspace{1.1mm}\par
{\small\textbf{Merke:} Jede Zeile der Hilfstabelle enthält die Primärschlüssel beider Tabellen als Fremdschlüssel. Aus der n:m-Beziehung werden zwei 1:n-Beziehungen: \texttt{schueler} \textbf{1--n} \texttt{teilnahme} \textbf{n--1} \texttt{ag}. Eine neue Anmeldung (Ben im Chor) ergänzt eine Zeile mit \texttt{schueler\_id = 2} und \texttt{ag\_id = 3}.}\par

\sectiontitle{3 · Hilfstabellen in InstaHub}
{\small\begin{tabularx}{\linewidth}{@{}>{\ttfamily\bfseries\color{Blue}}p{24mm}X@{}}
likes & \texttt{user\_id = 18, photo\_id = 14}: Benutzer 18 hat Foto 14 geliked. \texttt{photos.user\_id} nennt dagegen die Person, die das Foto hochgeladen hat.\\[1mm]
follows & \texttt{follower\_id} bezeichnet, wer folgt; \texttt{following\_id}, wem gefolgt wird. Beide Fremdschlüssel verweisen auf \texttt{users}. Gegenseitiges Folgen braucht zwei Zeilen.\end{tabularx}}\par\vspace{1mm}
\begin{tcolorbox}[colback=Mint,colframe=Blue,arc=2mm,boxrule=.8pt,left=3mm,right=3mm,top=1.7mm,bottom=1.7mm]
\textbf{Merksatz}\par\smallskip
Bei 1:1 steht der Fremdschlüssel in einer der beiden Tabellen. Bei n:m werden die Zuordnungen in einer Hilfstabelle mit zwei Fremdschlüsseln gespeichert.
\end{tcolorbox}

\vfill{\color{Line}\rule{\linewidth}{.5pt}}\par\vspace{1.5mm}
\begin{tabularx}{\linewidth}{@{}Xr@{}}{\small\color{Muted}Klasse 10 · Datenbanken} & {\small\color{Muted}Sicherungsblatt · Aufgabe 8}\end{tabularx}
\end{document}
